% SGA 6, Exposición VIII: fragmento español alineado con la fuente.
% Alcance exacto: sección 3 completa, fr.tex, líneas 11452--11648.
% SHA-256 de la fuente francesa completa:
% 0D2158B9FE3B466309A92165727A0B191422E05FBAA0DACD5806E8322C54BA1A.
% Termina inmediatamente antes de la sección 4, sin cierre sintético.
% Las superficies matemáticas se conservan diplomáticamente; las anomalías
% detectadas se notifican por separado para su resolución canónica.

\subsection*{3. Teorema de Riemann--Roch: enunciado}

\textbf{3.1.} En el número 3, $Y$ es un esquema cuasicompacto que posee un haz
amplio.

Sea $f:X\to Y$ un morfismo proyectivo, de intersección completa (1.1); se
supone además que la dimensión relativa virtual de $f$ (1.9) es constante sobre
$X$. Como $f$ es proyectivo, existe una inmersión cerrada $X\to P$, donde $P$
es un fibrado proyectivo sobre $Y$, y, como $Y$ posee un haz amplio, se puede
suponer que $P=\mathbb P^r_Y$. Se tiene, pues, una factorización
\[
\begin{tikzcd}[column sep=large,row sep=large]
X \arrow[r,"i"] \arrow[dr,"f"'] & \mathbb P^r_Y \arrow[d,"g"]\\
& Y.
\end{tikzcd}
\]
Como $f$ es de intersección completa, la inmersión $i$ es regular (1.2).

Por 1.7, $f$ es un morfismo perfecto. Además, $Y$ posee un haz amplio y, por
consiguiente, también $X$ y $P$. Se puede definir entonces un homomorfismo
aditivo
\[
f_*:K^{\bullet}(X)\longrightarrow K^{\bullet}(Y)
\]
(IV 2.12), pues los $K^{\bullet}$ definidos en términos de haces localmente
libres de tipo finito o de complejos perfectos son idénticos. También se pueden
definir homomorfismos aditivos
\[
i_*:K^{\bullet}(X)\longrightarrow K^{\bullet}(P);
\qquad
 g_*:K^{\bullet}(P)\longrightarrow K^{\bullet}(Y),
\]
y se tiene $f_*=g_*\circ i_*$.

Sea $d$ la dimensión relativa virtual de $X$ sobre $Y$. La codimensión de $i$
es entonces $r-d$. Por VII 4.6, el homomorfismo $i_*$ da, después de tensorizar
por $\mathbb Q$, inclusiones
\[
i_*\bigl(\Fil^kK^{\bullet}(X)_{\mathbb Q}\bigr)
\subset
\Fil^{k+r-d}K^{\bullet}(P)_{\mathbb Q}.
\]
Por VI 5.8, el homomorfismo $g_*$ da inclusiones
\[
g_*\bigl(\Fil^{k'}K^{\bullet}(P)\bigr)
\subset
\Fil^{k'-r}K^{\bullet}(Y).
\]
Componiendo ambas, se obtiene:

\textbf{Proposición 3.2.} Sea $f:X\to Y$ un morfismo proyectivo, de
intersección completa, cuya dimensión relativa virtual $d$ es constante. Se
supone que $Y$ es cuasicompacto y admite un haz amplio. Entonces, para todo
$k\in\mathbb Z$, se tiene
\[
f_*\bigl(\Fil^kK^{\bullet}(X)_{\mathbb Q}\bigr)
\subset
\Fil^{k-d}K^{\bullet}(Y)_{\mathbb Q}.
\]
Se obtiene, pues, en particular, un homomorfismo aditivo
\[
f_{\mathrm{gr}}:\Gr K^{\bullet}(X)_{\mathbb Q}\longrightarrow
\Gr K^{\bullet}(Y)_{\mathbb Q}
\]
que envía $\Gr^kK^{\bullet}(X)_{\mathbb Q}$ en
$\Gr^{k-d}K^{\bullet}(Y)_{\mathbb Q}$.

\textbf{3.3.} Recordemos que en (V 6.1) hemos definido un funtor
$A\longmapsto \operatorname{Ch}(A)$ sobre la categoría de las $K$-álgebras
graduadas, con valores en la categoría de las $K$-$\lambda$-álgebras
aumentadas, siendo $K$ un anillo binomial (V 2.7). Si $A$ es una $K$-álgebra
graduada, se tiene
\[
\operatorname{Ch}(A)=K\times(1+\widehat A^+).
\]
Para toda $K$-$\lambda$-álgebra aumentada $\Lambda$, hemos definido un
$K$-$\lambda$-homomorfismo, denominado «~clase de Chern completada~» (V 6.7),
\[
\widetilde c=(\varepsilon,c):\Lambda\longrightarrow\operatorname{Ch}(\Gr_{\lambda}\Lambda).
\]
Por último, hemos definido un homomorfismo de anillos, funtorial en $A$ y
denominado carácter de Chern (V 6.4),
\[
\ch:\operatorname{Ch}(A)\longrightarrow K\oplus\widehat A^+_{\mathbb Q},
\]
y un endomorfismo de grupos, denominado carácter de Todd (V 1.16),
\[
\Todd:1+\widehat A^+\longrightarrow1+\widehat A^+.
\]
Aplicaremos estas definiciones a los $\lambda$-anillos aumentados
$K^{\bullet}(X)$ y $K^{\bullet}(Y)$, pues los anillos de aumentación
$H^0(X,\mathbb Z)$ y $H^0(Y,\mathbb Z)$ son binomiales (V 2.8). Para simplificar
las notaciones, pondremos, para todo esquema $Z$ y todo
$z\in K^{\bullet}(Z)$,
\[
\ch(z)=\ch(\widetilde c(z));
\qquad
\Todd(z)=\Todd(c(z)).
\]
Para toda álgebra graduada $A$, consideraremos el grupo $1+\widehat A^+$ como
contenido en $\widehat A$. Finalmente, si se tiene un homomorfismo $f:A\to B$,
aditivo y compatible con las graduaciones salvo por un desplazamiento, se
denotará también por $f$ el homomorfismo $\widehat A\to\widehat B$ que se deduce
de él.

\textbf{3.4.} Sea $X$ un esquema. Para todo $\mathcal O_X$-módulo localmente
libre de tipo finito $E$, se denota por $\check E$ la clase en $K^{\bullet}(X)$
del dual $\check E$ de $E$. Está claro que la función que a $E$ asocia
$\check E$ es una función aditiva sobre la categoría de los
$\mathcal O_X$-módulos localmente libres de tipo finito y, por consiguiente,
define un endomorfismo de $K^{\bullet}(X)$, que se denotará por
$x\rightsquigarrow\check x$.

\textbf{Proposición 3.5.} Sea $X$ un esquema cuasicompacto. Entonces, para todo
$x\in K^{\bullet}(X)$ y todo $i\geq0$, se tiene
\[
c^i(\check x)=(-1)^i c^i(x).
\]
La relación que se debe demostrar también puede escribirse
\[
c_t(\check x)=c_{-t}(x),
\]
poniendo $c_t(x)=\sum_{i\geq0}c^i(x)t^i$.
Por aditividad, se puede suponer que $x$ es la clase de un
$\mathcal O_X$-módulo localmente libre de tipo finito $E$. Además,
reemplazando $X$ por el esquema de banderas de $E$ si es necesario, lo cual da
una inyección sobre los $\Gr K^{\bullet}$ (VI 5.5), se puede suponer que $E$ es
un haz invertible. Se tiene entonces $\check E=E^{-1}$; por tanto,
$\check x=x^{-1}$.

Los $c^i(\check x)$ y $c^j(x)$ son nulos para $i,j\geq2$. Basta, pues,
demostrar que
\[
c^1(x^{-1})=-c^1(x).
\]
Ahora bien, $c^1(x)$ es la clase en $\Gr^kK^{\bullet}(X)$ de $x-1$, y
$c^1(x^{-1})$ la de $x^{-1}-1$. La igualdad
\[
x^{-1}-1=1-x+(x-1)(1-x^{-1})
\]
da
\[
x^{-1}-1\equiv1-x\pmod{\Fil^2K^{\bullet}(X)};
\]
de donde se deduce el resultado.

\textbf{Teorema 3.6 (Riemann--Roch).} Sea $f:X\to Y$ un morfismo proyectivo de
intersección completa (1.1), cuya dimensión relativa virtual $d$ es constante
(1.9). Se supone que $Y$ es cuasicompacto y admite un haz amplio. Con las
notaciones de 3.3, para todo $x\in K^{\bullet}(X)$ se tiene entonces la relación
\[
\ch(f_*(x))=f_{\mathrm{gr}}\bigl(\ch(x)\cdot\Todd(\check T_f)\bigr),
\]
donde $\check T_f$ denota el elemento dual de la clase del complejo cotangente
relativo $L^{X/Y}_\bullet$ (2.5 y 2.3).

En otros términos, el diagrama
\[
\begin{tikzcd}[column sep=huge,row sep=large]
K^{\bullet}(X) \arrow[r,"\ch\circ\widetilde c"] \arrow[d,"f_*"'] & \Gr K^{\bullet}(X)_{\mathbb Q} \arrow[d,"x\mapsto f_{\mathrm{gr}}(x\cdot\Todd(\check T_f))"]\\
K^{\bullet}(Y) \arrow[r,"\ch\circ\widetilde c"'] & \Gr K^{\bullet}(Y)_{\mathbb Q}
\end{tikzcd}
\]
es conmutativo.



\subsection*{Continuación de la demostración del teorema 3.6}

La demostración del teorema 3.6 ocupará el resto de esta exposición. Demos
algunos preliminares.

\textbf{Proposición 3.7.} Cualesquiera que sean
$x\in \Gr K^{\bullet}(X)_{\mathbb Q}$ e
$y\in \Gr K^{\bullet}(Y)_{\mathbb Q}$, se tiene
\[
 f_{\mathrm{gr}}\bigl(x\cdot f^{\mathrm{gr}}(y)\bigr)
 =
 f_{\mathrm{gr}}(x)\,y .
\]

Por aditividad, se puede suponer que \(x\) e \(y\) son homogéneos; sea
$x\in \Gr^iK^{\bullet}(X)_{\mathbb Q}$,
$y\in \Gr^jK^{\bullet}(Y)_{\mathbb Q}$.
Sean $x'\in \Fil^iK^{\bullet}(X)_{\mathbb Q}$,
$y'\in \Fil^jK^{\bullet}(Y)_{\mathbb Q}$ elementos que levantan \(x\) e
\(y\). Entonces la fórmula de proyección da
\[
 f_*\bigl(x'\cdot f^*(y')\bigr)=f_*(x')\,y'.
\]
Por 3.2, esta igualdad se escribe en
\(\Fil^{i+j-d}K^{\bullet}(Y)_{\mathbb Q}\). De ella se deduce una igualdad en
\(\Gr^{i+j-d}K^{\bullet}(Y)_{\mathbb Q}\), y se comprueba inmediatamente que
es precisamente la igualdad buscada.

\textbf{Lema 3.8.} Sean \(f:X\to Y\) y \(g:Y\to Z\) dos morfismos que
satisfacen las hipótesis de 3.6. Entonces \(g\circ f\) también las satisface.
Si \(x\) es un elemento de \(K^{\bullet}(X)\), y si 3.6 se verifica para
\(f\) en el elemento \(x\), y para \(g\) en el elemento \(f_*(x)\), entonces
se verifica para \(g\circ f\) en el elemento \(x\).

Por hipótesis, se tiene
\begin{equation}\tag{3.8.1}
 \ch(f_*(x))=
 f_{\mathrm{gr}}\bigl(\ch(x)\cdot\Todd(\check T_f)\bigr),
\end{equation}
y
\begin{equation}\tag{3.8.2}
 \ch(g_*(f_*(x)))=
 g_{\mathrm{gr}}\bigl(\ch(f_*(x))\cdot\Todd(\check T_g)\bigr).
\end{equation}
Combinando (3.8.1) y (3.8.2), se obtiene
\[
\ch(g_*(f_*(x)))=
g_{\mathrm{gr}}\!\left(
 f_{\mathrm{gr}}\bigl(\ch(x)\Todd(\check T_f)\bigr)\cdot \Todd(\check T_g)
\right).
\]
Aplicando 3.7, se obtiene
\[
\ch(g_*(f_*(x)))=
g_{\mathrm{gr}}\!\left(
 f_{\mathrm{gr}}\bigl(\ch(x)\cdot\Todd(\check T_f)\cdot
 f^{\mathrm{gr}}(\Todd(\check T_g))\bigr)
\right).
\]
La funtorialidad de \(c\) y de \(\Todd\) implica
\[
 f^{\mathrm{gr}}(\Todd(\check T_g))=\Todd(f^*(\check T_g)).
\]
Por 2.7, se tiene
\[
 \check T_{g\circ f}=\check T_f+f^*(\check T_g).
\]
La multiplicatividad del carácter de Todd implica, por tanto,
\[
 \Todd(\check T_f)\cdot f^{\mathrm{gr}}(\Todd(\check T_g))
 =
 \Todd(\check T_{g\circ f}).
\]
Se obtiene finalmente
\[
 \ch(g_*(f_*(x)))=
 g_{\mathrm{gr}}\bigl(f_{\mathrm{gr}}\bigl(\ch(x)\cdot\Todd(\check T_{g\circ f})\bigr)\bigr),
\]
que es el resultado anunciado.

\textbf{3.9.} Si se aplica 3.8 a la factorización de \(f\) dada en 3.1, se ve
que, para demostrar el teorema de Riemann--Roch, basta demostrarlo en el caso
de una inmersión regular de codimensión constante y en el caso del morfismo
estructural $\mathbb P^r_Y\longrightarrow Y$.
Además, se puede suponer que la inmersión \(i\) satisface la condición de VII
4.1. Sea, en efecto, \(s\) la sección de \(\mathbb P^2_{\mathbb P^r_Y}\)
introducida en VII 4.2. Entonces \(s\circ i\) es una inmersión regular que
satisface la condición de VII 4.1 (VII 4.2) y, por 3.8, basta demostrar 3.6 para
\(s\circ i\), \(\mathbb P^2_{\mathbb P^r_Y}\longrightarrow\mathbb P^r_Y\) y
\(\mathbb P^r_Y\longrightarrow Y\).

Se demostrará, pues, 3.6 demostrando los dos casos particulares siguientes:
\begin{enumerate}[label=\roman*)]
\item \(f\) es una inmersión regular de codimensión constante que satisface la
condición de VII 4.1;
\item \(f\) es el morfismo estructural \(\mathbb P^r_Y\to Y\).
\end{enumerate}
